\chapter[Recommenders I: Classical Methods]{Recommenders I: Classical Methods + Learning-to-Rank}

\section{Learning objectives}
\begin{itemize}
  \item Explain the business role of recommendation and ranking systems in e-commerce discovery.
  \item Distinguish memory-based collaborative filtering, matrix factorization, and learning-to-rank approaches.
  \item Understand the difference between explicit and implicit feedback, and why e-commerce is usually dominated by implicit signals.
  \item Evaluate recommendation systems using offline ranking metrics while recognizing the limits of offline evaluation.
\end{itemize}

\section{Why recommenders matter in e-commerce}
Recommendation systems are central to modern e-commerce because large catalogs create a matching problem.
Customers rarely inspect all available items.
The platform therefore needs mechanisms for deciding which products to surface, in what order, and in which context.

Recommendation is not only a technical personalization problem.
It is also a business control layer for:
\begin{itemize}
  \item product discovery,
  \item cross-sell and upsell,
  \item session continuation,
  \item re-engagement,
  \item merchandising efficiency,
  \item and marketplace supply exposure.
\end{itemize}

This chapter focuses on classical recommenders and ranking methods.
These methods remain important because they are interpretable, computationally efficient, and often strong baselines even when deep-learning systems are eventually adopted.

\section{Recommendation surfaces and tasks}
Not all recommender problems are identical.
A commerce platform usually contains several recommendation surfaces with different goals.

\subsection{Common recommendation tasks}
\begin{itemize}
  \item \textbf{Homepage recommendations:} highlight likely-interest products for returning or anonymous users.
  \item \textbf{Product detail recommendations:} suggest similar, complementary, or substitute items.
  \item \textbf{Cart recommendations:} support bundle expansion and accessory attachment.
  \item \textbf{Email or push recommendations:} re-engage users outside the session.
  \item \textbf{Search ranking and re-ranking:} order retrieved items for a query.
\end{itemize}

\subsection{Retrieval versus ranking}
It is useful to distinguish two subproblems:
\begin{itemize}
  \item \textbf{Candidate generation:} find a manageable subset of items from a very large catalog.
  \item \textbf{Ranking or re-ranking:} order the candidates for the current context.
\end{itemize}

This chapter emphasizes ranking-oriented thinking and classical candidate selection methods.
Chapter 8 will extend this to deep retrieval and embedding-heavy systems.

\section{Feedback in commerce recommendation}
Recommendation systems learn from user behavior.
The nature of that behavior strongly affects modeling choices.

\subsection{Explicit feedback}
Explicit feedback is directly provided by users, such as:
\begin{itemize}
  \item star ratings,
  \item likes or dislikes,
  \item structured review scores,
  \item preference selections.
\end{itemize}

Explicit feedback is valuable because it is interpretable, but in many commerce settings it is sparse.
Most customers do not rate enough items to support the full recommendation problem.

\subsection{Implicit feedback}
Implicit feedback is inferred from behavior:
\begin{itemize}
  \item clicks,
  \item product views,
  \item add-to-cart events,
  \item purchases,
  \item dwell time,
  \item repeat visits,
  \item returns or cancellations as negative indicators.
\end{itemize}

E-commerce is largely an implicit-feedback domain.
This makes modeling harder because a missing interaction does not necessarily mean dislike.
It may simply mean the user never saw the item.

\subsection{Signal strength}
Not all implicit signals are equally informative.
For example:
\begin{itemize}
  \item a purchase usually signals stronger intent than a click,
  \item an add-to-cart event may indicate interest but not final preference,
  \item a return may weaken confidence in the original purchase signal,
  \item repeated category browsing may indicate exploration rather than commitment.
\end{itemize}

Feature and label design should therefore reflect business meaning, not just event counts.

\section{Memory-based collaborative filtering}
Classical collaborative filtering starts from the idea that similar users may like similar items, and similar items may be preferred by similar users.

\subsection{User-based collaborative filtering}
In user-based collaborative filtering, the system identifies users whose historical behavior resembles the target user and recommends items those similar users engaged with.
Conceptually:
\begin{itemize}
  \item define a user--item interaction matrix $R=[r_{ui}]$,
  \item compute user-user similarity,
  \item aggregate neighbor preferences for unseen items.
\end{itemize}

This approach is intuitive but often scales poorly for large consumer systems with many users and rapidly changing behavior.

\subsection{Item-based collaborative filtering}
Item-based collaborative filtering compares items rather than users.
If two items are frequently co-viewed, co-purchased, or similarly interacted with, one can recommend one when the other is observed.

Item-based methods are often more stable in commerce because items change more slowly than user sessions.
They work well for:
\begin{itemize}
  \item ``similar items'',
  \item ``customers also bought'',
  \item and complementary recommendation widgets.
\end{itemize}

\subsection{Similarity measures}
Common similarity measures include cosine similarity, Pearson correlation, and co-occurrence-based heuristics.
The choice depends on data type and normalization assumptions.
Architecturally, what matters is that similarity calculations must be based on well-defined interaction windows and item identities.

\section{Matrix factorization}
Memory-based methods are useful, but they struggle with sparsity and scale.
Matrix factorization addresses this by learning low-dimensional latent representations for users and items.

\subsection{Basic idea}
Given an interaction matrix $R$, matrix factorization approximates it as:
\[
\hat{r}_{ui}=\mathbf{p}_u^\top \mathbf{q}_i
\]
where $\mathbf{p}_u$ is a latent vector for user $u$ and $\mathbf{q}_i$ is a latent vector for item $i$.

These latent factors capture hidden structure such as taste patterns, style clusters, or price sensitivity without requiring manually defined categories.

\subsection{Why factorization works well}
Matrix factorization often performs strongly because:
\begin{itemize}
  \item it compresses sparse interactions into learnable structure,
  \item it generalizes beyond direct co-occurrence,
  \item it supports efficient ranking once vectors are learned,
  \item it forms a bridge toward later embedding-based methods.
\end{itemize}

\subsection{Limitations}
Despite its strengths, factorization has limits:
\begin{itemize}
  \item it handles cold-start users and items poorly without additional features,
  \item it may ignore rich session context,
  \item it can struggle when preferences shift quickly,
  \item it usually needs careful treatment of implicit feedback and negative sampling.
\end{itemize}

\section{Implicit-feedback recommendation}
Most commerce recommenders do not predict star ratings.
Instead, they infer preference from implicit interactions.
This changes both objective design and evaluation.

\subsection{Positive-only observations}
In implicit-feedback settings, observed interactions are often treated as positive evidence, while unobserved interactions are ambiguous.
Training usually requires assumptions such as:
\begin{itemize}
  \item sampling unobserved items as negatives,
  \item weighting confidence by interaction type,
  \item or optimizing pairwise ranking losses instead of rating prediction.
\end{itemize}

\subsection{Confidence weighting}
Different interactions can be assigned different strengths.
For example:
\begin{itemize}
  \item purchase $>$ add-to-cart $>$ click $>$ impression,
  \item repeat purchase may strengthen confidence,
  \item return or refund events may weaken confidence.
\end{itemize}

This is not just a modeling choice.
It is a business-rule decision and should align with how the organization interprets engagement quality.

\section{Cold start and sparsity}
Cold start is one of the defining challenges of recommendation systems.
It occurs when the system has insufficient data about a new user, a new item, or a new marketplace participant.

\subsection{User cold start}
New users may have little or no history.
Common mitigation strategies include:
\begin{itemize}
  \item popularity-based defaults,
  \item category-entry onboarding questions,
  \item contextual recommendations from landing page or query behavior,
  \item session-based short-term signals.
\end{itemize}

\subsection{Item cold start}
New items may not yet have interaction data.
Mitigation strategies include:
\begin{itemize}
  \item using metadata such as category, brand, and price range,
  \item manual merchandising boosts,
  \item seller and catalog-quality priors,
  \item controlled exploration to gather initial interactions.
\end{itemize}

\subsection{Marketplace and fairness considerations}
In marketplaces, recommendation systems also influence seller exposure.
If the system only amplifies already popular items, new or smaller sellers may never gather enough interaction data to compete.
This creates a business and governance concern, not just a model-performance issue.

\section{Learning-to-rank}
Recommendation is often framed as a ranking problem rather than a rating-prediction problem.
Learning-to-rank (LTR) methods optimize the order of items shown to users for a given context.

\subsection{Pointwise, pairwise, and listwise views}
LTR methods are often grouped into three families:
\begin{itemize}
  \item \textbf{Pointwise:} predict a score or label for each item independently.
  \item \textbf{Pairwise:} learn which of two items should be ranked higher.
  \item \textbf{Listwise:} optimize over the ordering of whole ranked lists.
\end{itemize}

\subsection{Why LTR is important in commerce}
E-commerce surfaces almost always have limited slots.
The platform is not merely predicting whether an item is relevant; it is deciding which few items to prioritize.
LTR is therefore especially appropriate for:
\begin{itemize}
  \item search result ordering,
  \item recommendation carousel ranking,
  \item category page re-ranking,
  \item email content selection.
\end{itemize}

\subsection{Feature design for LTR}
Classical ranking models often use hand-engineered or structured features such as:
\begin{itemize}
  \item user affinity to category or brand,
  \item historical click-through rate,
  \item purchase propensity,
  \item price distance from user norms,
  \item popularity and freshness,
  \item inventory availability,
  \item promotion eligibility,
  \item query-item lexical match for search contexts.
\end{itemize}

This makes LTR highly practical because structured business features are often already available in commerce systems.

\section{Offline evaluation}
Recommendation systems need offline evaluation for rapid iteration, model comparison, and failure analysis before live deployment.

\subsection{Common ranking metrics}
Frequently used metrics include:
\begin{itemize}
  \item precision@$k$,
  \item recall@$k$,
  \item mean average precision (MAP@$k$),
  \item normalized discounted cumulative gain (NDCG@$k$),
  \item mean reciprocal rank (MRR) in some search-style tasks.
\end{itemize}

\subsection{What these metrics measure}
\begin{itemize}
  \item \textbf{Precision@$k$:} how many of the top-$k$ items are relevant.
  \item \textbf{Recall@$k$:} how much of the relevant set is recovered in the top-$k$.
  \item \textbf{NDCG@$k$:} rewards placing more relevant items higher in the ranked list.
\end{itemize}

These metrics are useful because they reflect ranking quality rather than raw score accuracy.

\subsection{Limits of offline evaluation}
Offline metrics are not enough on their own.
They cannot fully capture:
\begin{itemize}
  \item causal business impact,
  \item position bias and feedback loops,
  \item changing user behavior under new experiences,
  \item interference with other surfaces such as promotions or search.
\end{itemize}

This is why offline evaluation should be seen as a filtering step, not the final decision criterion.

\section{Online evaluation and business impact}
Ultimately, recommender systems exist to improve user and business outcomes in production.
Typical online measures include:
\begin{itemize}
  \item click-through rate,
  \item add-to-cart rate,
  \item conversion rate,
  \item average order value,
  \item revenue per session,
  \item repeat engagement,
  \item return rate or dissatisfaction signals.
\end{itemize}

However, online KPI interpretation must be careful.
A recommender can increase clicks while harming profitability or reducing diversity.
Guardrail metrics are therefore important.

\section{Biases and practical risks}
Recommendation systems are affected by several structural biases.

\subsection{Popularity bias}
Popular items accumulate more exposure and therefore more interactions, reinforcing their future dominance.
This can suppress niche discovery and reduce long-tail coverage.

\subsection{Selection bias}
Models only observe interactions for items that were shown.
Unseen items do not generate evidence, so naive training pipelines may overestimate the quality of historically exposed items.

\subsection{Position bias}
Higher-ranked items receive more clicks partly because they are higher ranked, not only because they are more relevant.
This complicates interpretation of click data.

\subsection{Business-constraint tension}
The ``best'' recommendation by predicted engagement may conflict with:
\begin{itemize}
  \item margin objectives,
  \item inventory constraints,
  \item seller-fairness rules,
  \item compliance or policy requirements,
  \item user-experience diversity needs.
\end{itemize}

Recommendation is therefore not a pure prediction task.
It is a constrained decision problem embedded in business policy.

\section{A practical recommender architecture}
Even for classical methods, recommender systems are not just models.
They require an end-to-end architecture that includes:
\begin{itemize}
  \item event collection and identity linking,
  \item feature generation,
  \item offline training,
  \item candidate generation,
  \item ranking logic,
  \item serving interfaces,
  \item monitoring and experimentation.
\end{itemize}

This chapter remains focused on the modeling layer, but students should keep the surrounding system in mind.
The model only creates value when the architecture can serve it reliably and measure its effect.

\section{Hands-on lab: a baseline recommender}
The hands-on goal for this chapter is to build and evaluate a classical recommendation baseline.

\subsection{Lab scenario}
Students are given a synthetic or public interaction dataset containing user IDs, item IDs, timestamps, and event types such as view, add-to-cart, and purchase.

\subsection{Lab tasks}
\begin{enumerate}
  \item Construct a user--item interaction matrix using implicit feedback.
  \item Implement a simple popularity baseline and one collaborative filtering or matrix-factorization baseline.
  \item Evaluate the models using MAP@$K$ and NDCG@$K$.
  \item Analyze cold-start and sparsity cases explicitly.
  \item Propose one feature-based re-ranking rule that incorporates a business constraint such as inventory or diversity.
\end{enumerate}

\subsection{Expected learning outcome}
By the end of the lab, students should be able to explain not only which model performed better offline, but why the evaluation setup and business constraints matter.

\section{Managerial implications}
Classical recommendation methods remain valuable in practice because they provide:
\begin{itemize}
  \item strong baselines,
  \item fast iteration cycles,
  \item reasonable interpretability,
  \item lower infrastructure cost than many deep systems.
\end{itemize}

Organizations often make a mistake by jumping directly to deep-learning architectures without first establishing strong baselines, high-quality data, and trustworthy evaluation pipelines.
For many commerce surfaces, the real bottleneck is not model sophistication but poor signal design or unclear business objectives.

\section{Multiple-choice questions (MCQs)}
\begin{enumerate}
  \item In e-commerce recommendation, which statement best distinguishes \emph{implicit} from \emph{explicit} feedback?
  \begin{enumerate}
    \item Implicit feedback is always more accurate than explicit feedback.
    \item Explicit feedback comes directly from user judgments such as ratings, whereas implicit feedback is inferred from behaviors such as clicks or purchases.
    \item Explicit feedback can only be used in search ranking.
    \item Implicit feedback means the platform has no user identifiers.
  \end{enumerate}

  \item Item-based collaborative filtering is often attractive in commerce because:
  \begin{enumerate}
    \item items are usually more stable than users and item-item similarities can support scalable recommendation widgets.
    \item it eliminates cold start entirely.
    \item it requires no interaction data.
    \item it guarantees causal lift.
  \end{enumerate}

  \item In matrix factorization, the expression $\hat{r}_{ui}=\mathbf{p}_u^\top \mathbf{q}_i$ is intended to represent:
  \begin{enumerate}
    \item a network timeout policy.
    \item the predicted affinity between user $u$ and item $i$ using latent factors.
    \item the exact accounting value of an order.
    \item a lossless compression ratio for logs.
  \end{enumerate}

  \item Which metric is most appropriate for evaluating the quality of a top-$k$ ranked recommendation list offline?
  \begin{enumerate}
    \item NDCG@$k$
    \item Mean absolute error on shipping cost
    \item CPU utilization
    \item Database replication lag
  \end{enumerate}

  \item Why is learning-to-rank especially relevant in e-commerce?
  \begin{enumerate}
    \item Because commerce systems never need candidate generation.
    \item Because the platform must decide which limited set of items to place at the top of a ranked surface.
    \item Because it replaces online experimentation.
    \item Because it removes the need for features.
  \end{enumerate}
\end{enumerate}

\subsection*{Answer key}
\begin{enumerate}
  \item (b)
  \item (a)
  \item (b)
  \item (a)
  \item (b)
\end{enumerate}

\section{Exercises (short)}
\begin{enumerate}
  \item Compare user-based, item-based, and matrix-factorization approaches for a medium-sized commerce catalog. State one strength and one weakness of each.
  \item Define an implicit-feedback weighting scheme for view, add-to-cart, purchase, and return events. Explain the business reasoning behind your weights.
  \item Design a simple learning-to-rank feature set for a ``recommended for you'' widget. Include at least one user feature, one item feature, one context feature, and one business-rule feature.
\end{enumerate}

\section{Mini-case (Odoo-linked, vendor-neutral)}
\textbf{Scenario:} A company uses a headless storefront, Odoo for ERP, and a central data platform. The business wants recommendation widgets on the homepage, product pages, and cart page.

\textbf{Task:} Propose a classical recommender design:
\begin{itemize}
  \item Define which interaction signals from storefront and Odoo-derived outcomes should be used for training.
  \item Choose a baseline method (item-based CF, matrix factorization, or feature-based ranking) for each surface and justify the choice.
  \item Identify two business constraints the ranking layer should respect, such as inventory, margin, seller exposure, or return risk.
\end{itemize}
